Hij keek je blinkend aan
als daglicht wist te gaan

door de bonte ooglens.
Hij was een kleurig mens

volgens zijn enig oog.
Men zag de regenboog

glinsteren, een symptoom
van staar als een kerstboom.

De oude vagabond
was overal verwond.

Door een pijl uit een boog
verdween zijn linkeroog

en werd een litteken,
dat ver uit bleef steken

boven diepe rimpels.
Hij straalde iets simpels,

luchtig, maar robuust uit.
Zijn verweerde hoofdhuid

droeg een zwartgrijze baard.
Zijn nonchalante aard

leidde tot versloffing.
Dat bracht tot ontploffing

al zijn wilde haren.
Die had hij al jaren.

Als ruige jongeman
genoot hij volop van

wat het leven hem bood.
Onbezorgd om de dood

jaagde hij illegaal.
Hij stroopte zijn jachtmaal

in een jachtreservaat.
De elite werd kwaad

toen ze de effecten
op het wild ontdekten.

De buit uit zijn rooftocht
werd op de markt verkocht.

Die buit bestond later
uit elke verlater

van de lijfeigenschap.
Of een ander heerschap

dat vogelvrij rondzwierf
en niet in vrijheid stierf.

Begeerlijke vrouwen
wilden met hem trouwen

toen de jager nog jong
zijn toontje hoger zong.

De als verwilderde
boef afgeschilderde

contrasteerde met zoet.
Vooral het blauwe bloed

zag hem als bonte spijs.
Dat gebruikte hij wijs.

De façade, bevrijd
van verzorgde schoonheid,

liet zijn sluier vallen.
De jager mocht knallen.

Aldus werd de jager
meer een scharrelslager

dan een huwbare man.
Zijn misdadig elan

vergrootte de spanning.
Zijn rebelse houding

verwoordde hun eigen
sputterend stilzwijgen

over de dictatuur
van de adelcultuur.

Spijts zijn escapades
kwamen er geen schades,

bijgenaamd bastaarden,
die hij kon aanvaarden.

Hij woonde in een stad
waar hij uit de wind zat.
Door sociale onrust
ontstond daar de strijdlust

en de geloofsgekte
voor een nieuwe sekte.

Deze samenleving
slikte geen vervolging

door de nobiliteit.
Dat was het stadsbeleid

en dat wisten velen.
Zowaar criminelen

werd dekking geboden.
Ook de Griekse goden

op de berg Olympus
hoorden bij het vulgus

dat strijdt van tijd tot tijd.
Moraliteit ten spijt

werden ze wel vereerd.
Zonden geaccepteerd.

Dus waarom niet mensen
hetzelfde toewensen?

Zijn vader was zeer rijk
geworden door een dijk

aan zee aan te leggen
en land droog te leggen.

Door het recht van opstrek
claimde hij het als stek.

Ook bij de adelstand
vond hij meer boerenland.

Hij fikste een grondhuur
tegen een lange duur.

Zo bracht de huurprijsstop
door inflatie meer op,

omdat de winst toenam,
daar de oogst hem toekwam.

Op land in Groningen
verrezen woningen

voor zijn boerenknechten.
Door die beklemrechten

werd hij een herenboer.
De jager had een broer,

die vijf jaar jonger was
en niet zo’n wildebras,

maar altijd extreem moe.
Hij deed zijn ogen toe

dikwijls zelfs overdag.
Met slapheid als beklag

kreeg je bij de jager
snel het stempel 'klager'.

Zijn broertje viel vaak flauw
en verslikte zich gauw.

Zijn darmen prikkelden
als beiden smikkelden.

“Schijt jij nu je broek vol?”
De jager rook een drol.

Zijn broertje schudde nee,
maar keek ineens tevree

toen hij bleek te drukken.
“Het lijkt toch te lukken!”,
zei zijn broertje voldaan.
Hij had iets knaps gedaan,

waarmee hij vast scoorde.
De jongeling hoorde

zijn vader vaak roepen:
“Die vlijtige poepen!

Prachtig! Dat kan slechter!”
Vader verwees echter

naar seizoenarbeiders.
Het waren bevrijders

van de belegering
door werk uit turfwinning.

Ze kwamen uit Duitsland
en reisden naar Holland.

Vader legde eerloos
alswel meedogenloos

de jager en zijn broer,
opzichtelijk een loer.

Hij verried zijn zonen
om zich te verschonen

van een zedenmisdaad.
Moeder steunde zijn daad

en hielp zijn deceptie.
Met zo’n rot familie

zijn geen vetes nodig.
Het is overbodig

te melden dat de broers
snel verlegden hun koers

en geen boer meer werden.
De lachende derde

na de broers exodus,
was hun enige zus.

Zijn kinderwens ten spijt,
gaf zijn mannelijkheid

hem nog geen nageslacht.
Zijn kleine broertje bracht

daarom trots de tijding
dat hij in verwachting

was geraakt met zijn vrouw.
Die pesterij was flauw,

maar werd goed ontvangen.
Hun beider verlangen

om vader te worden
nam vlot alle horden

over zijn jaloezie.
Het was zijn empathie,

die hem oprecht heugde
de vaderschapsvreugde.

“Als ik geen nageslacht
op tijd heb voortgebracht

dan sterft mijn bloedlijn uit!”
maalde de jager luid.

Eerst was het de liefde
die hem intens griefde.

Het ongeboren kind
dat hij al had bemind.

Later was het zijn angst
te blijven zonder vangst,

zonder ontvangenis.
Deze gevangenis

bouwde een gouden kooi
door een gehaaid pleidooi

over mannelijkheid
en zijn sterfelijkheid

tegen de zinloosheid.
Generativiteit

werd een noodwendigheid
voor continuïteit.

Hij was geen voorste haan,
zoals ooit Genghis Khan

overal zijn genen
tussen vrouwenbenen

vruchtbaar kon verspreiden.
Hem wist te verblijden

als hij zijn genenprint
al in een enkel kind

vruchtbaar kon nalaten.
Hij zou niet nalaten

zijn zorgtaak te nemen
ondanks
drankproblemen.

De hoop op het goede,
die eerst geloof voedde
ondanks onvruchtbaarheid,
raakte de jager kwijt.

De toekomst beklemde.
Die begrensdheid remde

zijn voortplantingsdriften.
Hoewel de geschriften,

als hij die afspeurde,
hem soms weer opbeurden.

“Bij God is mogelijk
wat is onmogelijk

bij de mens”, zei Jezus
ooit volgens Matteüs.

De jager kon de hoop,
die hem steeds weer bekroop,

niet verwerkelijken.
Niet verwezenlijken.

Of was Emmy’s klaproos
wellicht Pandora’s doos?
